\begin{reviewcomment}{3}
	The policy language is still quite limited.
	The system focuses on fixed-length Boolean RPN programs, while many realistic healthcare consent policies involve temporal constraints, revocation, obligations, delegation, purpose limitation, institutional approval, and evolving consent.
	Some of these are treated as external predicate bits, but this moves complexity outside the proven statement.
	The paper should better explain what kinds of real consent policies can actually be supported end-to-end.
\end{reviewcomment}

\begin{reviewresponse}{3}
	We agree that ``external predicate bit'' left the guarantee unclear, and the revision answers the question in two steps: the relation no longer accepts unconstrained predicate bits, and the Design section now states feature by feature what the prototype establishes at authorization time.

	In the revised relation every predicate bit is derived in-circuit from a witness bound to committed parameters: ontology-scoped conditions (disease, study design, ancestry, geography, clinical term) are Merkle-membership proofs against a provider-committed whitelist root; user-, project-, and institution-specific restrictions are membership proofs of the holder pseudonym in provider-committed allowlists; and the profile-backed data-use terms are derived from the claims of an EdDSA-Poseidon credential that binds the policy commitment, the requester pseudonym, the resource, the claims, a validity interval, and a credential identifier.
	The policy's own validity window is committed and checked against the public request epoch, which the contract requires to equal the current hour.

	A new table in the Design section classifies each consent feature as decided in-circuit, decided on-chain by the access-control contract from public state, committed but not decided at request time, or not provided. The Conclusion lists the unsupported features among its limitations, and the Discussion suggests future work to extend the class further (a governed issuer registry with in-circuit membership and revocation witnesses, and holder-controlled release of claims) without changing the interpreter or the public interface, because new conditions enter through predicate slots and the fixed-format parameter commitment.


	\RevisionLocation{Section ``TPS-ZKPI Design,'' subsection ``Policy Representation and Interpretation,'' new paragraph ``Scope of the supported policies'' and the new policy-support table; Section ``Conclusion,'' third paragraph.}
	\begin{revisiontext}
		The policy class is deliberately restricted, and the policy-support table states what the RC1 prototype establishes end-to-end.
		Categorical data-use terms, ontology-scoped conditions, requester allowlists, profile-backed terms, the policy validity window, and time-limited use are decided at authorization time, either inside the proved relation or by the ACM from public state.
		A publication moratorium is bound into the policy as an accepted obligation but cannot be decided at request time, negation is rejected by the DUO compiler because membership witnesses prove inclusion only, and purpose limitation, post-access duties, delegation, and credential revocation are left to the data access agreement and the issuer lifecycle.
		The prototype is therefore an access-time consent checker for bounded Boolean policies over standardized terms, not a general consent-language engine.
	\end{revisiontext}
\end{reviewresponse}
